% ============================================================
%  Planning Algorithms Benchmark — Introductory Presentation
%  Compile: pdflatex main.tex  (x2 for cross-refs)
%  Requires standard TeX Live: beamer, booktabs
% ============================================================
\documentclass[aspectratio=169, 11pt]{beamer}

\usetheme{Boadilla}
\usecolortheme{default}

\definecolor{Dark}{HTML}{1E293B}
\definecolor{Mid}{HTML}{475569}
\definecolor{Light}{HTML}{CBD5E1}
\definecolor{LeWMc}{HTML}{0891B2}
\definecolor{DINOc}{HTML}{7C3AED}
\definecolor{BG}{HTML}{F8FAFC}

\setbeamercolor{structure}{fg=Dark}
\setbeamercolor{frametitle}{bg=Dark, fg=white}
\setbeamercolor{title}{fg=white}
\setbeamercolor{subtitle}{fg=Light}
\setbeamercolor{author}{fg=Light}
\setbeamercolor{date}{fg=Light}
\setbeamercolor{palette primary}{bg=Dark, fg=white}
\setbeamercolor{palette secondary}{bg=Dark!80, fg=white}
\setbeamercolor{palette tertiary}{bg=Dark!60, fg=white}
\setbeamercolor{palette quaternary}{bg=Dark, fg=white}
\setbeamercolor{block title}{bg=Dark, fg=white}
\setbeamercolor{block body}{bg=BG}
\setbeamercolor{block title alerted}{bg=DINOc, fg=white}
\setbeamercolor{block body alerted}{bg=DINOc!8}
\setbeamercolor{block title example}{bg=LeWMc, fg=white}
\setbeamercolor{block body example}{bg=LeWMc!8}
\setbeamertemplate{navigation symbols}{}
\setbeamertemplate{itemize items}[circle]

\usepackage[T1]{fontenc}
\usepackage{lmodern}
\usepackage{booktabs}
\usepackage{xspace}

\newcommand{\lewm}{\textcolor{LeWMc}{\textbf{LeWM}}\xspace}
\newcommand{\dino}{\textcolor{DINOc}{\textbf{DINO-WM}}\xspace}

\title{Planning with Learned World Models}
\subtitle{A Benchmarking Study}
\author{Anurag Dhiman}
\date{\today}

% ============================================================
\begin{document}

% ── TITLE ────────────────────────────────────────────────────
{
\setbeamertemplate{footline}{}
\begin{frame}
  \titlepage
\end{frame}
}

% ── 1. RESEARCH PROBLEM ──────────────────────────────────────
\begin{frame}{Research Problem}

  \vspace{2em}
  \begin{center}
    \large
    Given the same learned world model and the same planning budget,\\[0.6em]
    which planning algorithm performs best\\[0.6em]
    under different task characteristics?
  \end{center}

\end{frame}

% ── 2. INITIAL QUESTIONS ─────────────────────────────────────
\begin{frame}{Initial Questions to Explore}

  \begin{enumerate}
    \setlength\itemsep{1.1em}

    \item Does \textbf{world model quality} determine planning success?\\
    {\small Can a good planner compensate for an inaccurate dynamics model?}

    \item Does the \textbf{compute budget} $B$ (model calls per step) limit performance?\\
    {\small Is the planner failing because the search is too shallow?}

    \item Which \textbf{planning algorithm} is most robust to model errors?\\
    {\small Classical search? Sampling-based? Model-based with learned value?}

    \item Do different \textbf{environments} expose different bottlenecks?\\
    {\small A small open grid vs.\ a multi-room maze — do they stress different things?}

  \end{enumerate}

\end{frame}

% ── 3. EXPERIMENTAL SETUP ────────────────────────────────────
\begin{frame}{Experimental Setup}

  \begin{columns}[T]

    \column{0.30\textwidth}
    \textbf{Environments}
    \begin{itemize}
      \item MiniGrid \textbf{Empty-8$\times$8}\\
            {\small Open grid, 253 max steps, sparse reward}
      \vspace{0.5em}
      \item MiniGrid \textbf{FourRooms}\\
            {\small Rooms and doorways, 97 max steps, sparse reward}
    \end{itemize}

    \column{0.33\textwidth}
    \textbf{World Models}
    \begin{itemize}
      \item \dino\\
            {\small Frozen DINOv2 encoder\\with learned dynamics head}
      \vspace{0.5em}
      \item \lewm\\
            {\small ViT-tiny encoder\\trained with JEPA-style objective}
    \end{itemize}

    \column{0.33\textwidth}
    \textbf{Planning Algorithms}
    \begin{itemize}\small
      \item Classical: BFS, UCS, A*, Greedy-BFS
      \item Sampling: MCTS, CEM, Beam Search, Random Shooting
      \item Dynamic programming: Value Iteration, Policy Iteration
      \item Model-based: Value-Guided A*, MuZero-style
    \end{itemize}
    \vspace{0.5em}
    {\small \textbf{Budget:} $B = 200$ model calls/step\\
    \textbf{Episodes:} 20 per configuration}

  \end{columns}

\end{frame}

% ── 4. STARTING RESULTS ──────────────────────────────────────
\begin{frame}{Starting Results: Success Rate at $B = 200$}

  \begin{columns}[T]

    \column{0.47\textwidth}
    \begin{block}{MiniGrid Empty-8$\times$8}
      \begin{itemize}
        \item \dino: \textbf{0\%} across all planners
        \item \lewm: up to \textbf{20\%} (CEM, MCTS+Value)
        \item Most planners: 0--5\% with both models
        \item Budget $B = 200$ is not the limiting factor here —\\
              search fits within budget when the model is accurate
      \end{itemize}
    \end{block}

    \column{0.47\textwidth}
    \begin{block}{MiniGrid FourRooms}
      \begin{itemize}
        \item \dino: up to \textbf{5\%} (MuZero-style only)
        \item \lewm: up to \textbf{10\%} (MCTS+Value, MuZero-style)
        \item All planners and both models perform similarly
        \item Even with perfect dynamics, planners reach $\leq 10\%$ —\\
              the search space is too large for $B = 200$
      \end{itemize}
    \end{block}

  \end{columns}

  \vspace{1em}
  \begin{center}
    Both models achieve low success rates overall.\\
    But the \emph{reason} for failure differs between the two environments.
  \end{center}

\end{frame}

% ── 5. WORLD MODEL QUALITY ───────────────────────────────────
\begin{frame}{World Model Quality: Rollout Stability}

  \begin{columns}[T]

    \column{0.47\textwidth}
    \begin{alertblock}{\dino — unstable rollouts}
      \begin{itemize}
        \item Latent L2 error at horizon $H{=}1$: \textbf{6.5}
        \item Latent L2 error at horizon $H{=}20$: \textbf{11{,}823}
        \item Growth: $\times$\textbf{1807} over 20 steps
        \item Frozen DINOv2 features drift catastrophically\\
              under iterative self-prediction
      \end{itemize}
    \end{alertblock}

    \column{0.47\textwidth}
    \begin{exampleblock}{\lewm — stable rollouts}
      \begin{itemize}
        \item Latent L2 error at horizon $H{=}1$: \textbf{19.3}
        \item Latent L2 error at horizon $H{=}20$: \textbf{34.4}
        \item Growth: $\times$\textbf{1.8} over 20 steps
        \item JEPA-style training keeps the latent space\\
              consistent across prediction steps
      \end{itemize}
    \end{exampleblock}

  \end{columns}

  \vspace{0.9em}
  \begin{center}
    \lewm has worse 1-step accuracy but is far more suitable for multi-step planning.\\
    \textbf{Training objective matters more than encoder capacity.}
  \end{center}

\end{frame}

% ── 6. CURRENT CONCLUSIONS ───────────────────────────────────
\begin{frame}{Current Conclusions}

  \begin{enumerate}
    \setlength\itemsep{1.0em}

    \item \textbf{Two distinct failure modes, one per environment}\\
    {\small Empty-8$\times$8: world model quality is the bottleneck ($B=200$ is sufficient but models are wrong).\\
    FourRooms: budget is the bottleneck (even with perfect dynamics, $B=200$ is not enough).}

    \item \textbf{Rollout stability matters more than 1-step accuracy}\\
    {\small \lewm grows error $\times$1.8 over 20 steps; \dino grows $\times$1807. \lewm is far more useful for planning despite starting from a higher 1-step error.}

    \item \textbf{Planner choice has less impact than model quality}\\
    {\small On Empty-8$\times$8, all learned-model planners cluster at $\leq$20\% SR. Improving the model would lift every planner simultaneously.}

  \end{enumerate}

\end{frame}

% ── 7. NEXT QUESTIONS ────────────────────────────────────────
\begin{frame}{Next Questions to Solve}

  \begin{enumerate}
    \setlength\itemsep{1.1em}

    \item \textbf{How much budget does FourRooms actually need?}\\
    {\small Running a sweep over $B \in \{200,\, 500,\, 1000,\, 2000,\, 5000\}$ with the oracle model.\\
    Preliminary result: BFS success rate jumps from 5\% to 15\% at $B = 500$.}

    \item \textbf{Can we improve \lewm's 1-step accuracy without losing stability?}\\
    {\small \lewm starts at L2 error 19.3 vs \dino's 6.5 at $H=1$. A better training recipe could close this gap while preserving multi-step consistency.}

    \item \textbf{Which planner handles model errors most gracefully?}\\
    {\small Value-Guided A* and CEM are the strongest so far. Do they degrade less as model quality decreases?}

    \item \textbf{Does the bottleneck shift with environment scale?}\\
    {\small At what state-space size does budget become the dominant bottleneck even on simple navigation tasks?}

  \end{enumerate}

\end{frame}

% ── CLOSING ──────────────────────────────────────────────────
{
\setbeamertemplate{footline}{}
\begin{frame}[plain]
  \begin{center}
    \vspace{3em}
    {\Large\bfseries\color{Dark} Thank you}\\[1.5em]
    {\color{Mid} Questions?}
  \end{center}
\end{frame}
}

\end{document}
